\documentclass[11pt]{article}
\usepackage[a4paper,margin=28mm]{geometry}
\usepackage{amsmath,amssymb,amsthm,mathtools,hyperref}
\newcommand{\ddc}{dd^c}
\newcommand{\NP}[1]{\langle #1\rangle}
\newtheorem{theorem}{Theorem}
\newtheorem{lemma}[theorem]{Lemma}
\theoremstyle{remark}\newtheorem{remark}[theorem]{Remark}
\title{An explicit countable divisorial criterion for ideal positivity\newline
on a genuinely nonclosed projective surface}
\author{}\date{}
\begin{document}\maketitle
\noindent\textbf{Checked scoped A--F result.} Full proof comparison and
statement-only reconstruction have been completed with Xean and
\texttt{gpt-6.1-sol}, \texttt{xhigh}, relative to the cited classical
foundations. This is not a self-contained certification of those
foundations. Original P06 remains open.
\begin{abstract}
We prove a divisorial-order criterion for positive lower singularity
volume on a pluriclosed projective surface. Infinitely many hyperplane
orders specify a nonanalytic ideal type. A positive smooth residual gives
an explicit common finite mass for every representative of the ideal
envelope's ordinary type. The criterion uses all modification orders,
rather than the unknown envelope mass. A concrete nonclosed example has
infinitely many divisors outside a pencil. The known countable-divisor
envelope mechanism and bounded pluriclosed mass invariance receive their
existing credit; the conclusion combines them with an entire-type
non-pluripolar mass computation. General ideal positivity and the
zero-volume implication programme remain open.
\end{abstract}

Use $\ddc=\frac{i}{2\pi}\partial\bar\partial$ and ordinary wedge powers.
All potentials are genuine canonical quasi-plurisubharmonic representatives.
Write $v\preceq u$ for $v\leq u+C$, and $v\sim u$ for mutual order.
The ideal $\mathcal I(tu)$ uses $|g|^2e^{-tu}$ at every real $t>0$.
Extended products retain the non-pluripolar finite-level definition with
smooth background corrections. Their actual mass is distinguished from
the infimum and supremum over a whole ordinary type.

\section{The explicit criterion}
Let $X=\mathbb P^2$, with the standard Fubini--Study form $\kappa$,
$\int_X\kappa^2=1$. Let $\theta$ be a smooth positive real $(1,1)$-form
with $\ddc\theta=0$; it need not be closed. Choose pairwise distinct
hyperplanes $H_j$, positive real coefficients $a_j$, and
\[
 A=\sum_{j\geq1}a_j<\infty,\qquad R=\theta-A\kappa>0.
\]
For a normalized defining linear section $s_j$, set
\[
 \ell_j=\log\frac{|s_j(Z)|^2}{\|Z\|^2}\leq0,
 \qquad u_0=\sum_{j\geq1}a_j\ell_j,
 \qquad u=u_0-\sup_Xu_0.
\]
We use the generic divisor coefficient normalized by
$\nu_E(\log|g|^2)=\operatorname{ord}_E(g)$.
For every prime divisor $E$ on every smooth compact modification
$\pi:Z\to X$, define
\begin{equation}\label{eq:profile}
 b(E)=\sum_{j\geq1}a_j\operatorname{ord}_E(\pi^*H_j).
\end{equation}
The convergence and exact interpretation of this sum are proved below.

\begin{theorem}\label{th:criterion}
Suppose a genuine $\theta$-psh potential $v$ has
$\nu_E(v\circ\pi)=b(E)$ for every divisor and modification in
\eqref{eq:profile}. Equivalently, require
$\mathcal I(tv)=\mathcal I(tu)$ for every real $t>0$.
Let
\[
 q=P_\theta[v]_{\mathcal I}
 =\left(\sup\{f\in\operatorname{PSH}(X,\theta):f\leq0,
             \mathcal I(tf)=\mathcal I(tv)\ (t>0)\}\right)^*.
\]
Then $q$ is genuine and $q\sim u$. Every genuine $\theta$-psh
$w\sim q$ has the same finite actual mass
\begin{equation}\label{eq:mass}
 \int_X\operatorname{MA}^{\mathrm{ext}}_\theta(w)=\int_XR^2>0.
\end{equation}
Consequently both lower and upper singularity volumes of $q$ equal
this number. The background is big and $X$ has BMP. Neither a raw
envelope mass nor an unknown envelope is used as a positivity test.
\end{theorem}

The family has infinitely many nonzero prime orders, so its ideal type
has no analytic presentation on the compact base. The all-modification
condition is an exact sufficient criterion, with no claim that a finite
selection of divisors suffices or that it is a general classification.

\section{The logarithmic series and its ideal envelope}
The exact Poincar\'e--Lelong formula gives
$\ddc\ell_j=[H_j]-\kappa$. Unitary invariance and radial integration
give $\int_X\ell_j\kappa^2=-3/2$. Thus the decreasing series has finite
$L^1$ norm, has a genuine canonical limit, and is not identically
$-\infty$. With $S=\sum_ja_j[H_j]$, distributional convergence gives
\[
 \ddc u=S-A\kappa,\qquad \theta+\ddc u=R+S>0.
\]
This proves strict admissibility and bigness using $R$ itself as a
positive Hermitian form. For BMP choose $C\kappa\geq h$ for any fixed
Hermitian $h$. A bounded $h$-psh potential is $C\kappa$-psh and
\[
 0\leq(h+\ddc f)^2\leq(C\kappa+\ddc f)^2.
\]
The bounded closed product on the right has total mass $C^2$ by
the Bedford--Taylor recursion and Stokes, giving the required upper bound.

Fix $\pi$ and $E$ in \eqref{eq:profile}, and choose a smooth
Gauduchon form $\eta>0$ on the compact surface $Z$, so
$\ddc\eta=0$ \cite[Theorem 2]{Gau}. For a tail of total coefficient
$A_K=\sum_{j>K}a_j$, write $r_K=\sum_{j>K}a_j\ell_j$.
Its canonical pullback is genuine: it is finite somewhere in the open
set where $\pi$ is an isomorphism. The closed positive tail current
$T_K=A_K\pi^*\kappa+\ddc(r_K\circ\pi)$ satisfies, by distributional
Stokes and Siu's divisorial subtraction,
\[
 \nu_E(r_K\circ\pi)\int_E\eta
 \leq\int_Z T_K\wedge\eta
 =A_K\int_Z\pi^*\kappa\wedge\eta.
\]
Thus the tail's generic coefficient is at most $C_EA_K$, with the
explicit finite constant $C_E=(\int_Z\pi^*\kappa\wedge\eta)/(\int_E\eta)$.
Finite additivity and $A_K\to0$ prove
$\nu_E(u\circ\pi)=b(E)$, including convergence of the sum.

The all-weight, all-modification Lelong characterization of ideal
equivalence applies in a common sufficiently large closed background
$C\kappa$, not in a fictitious line-bundle interpretation of $\theta$.
For the implication from orders to ideals, use
\cite[Theorem 2.13 and Proposition 2.30]{DX}: the hypothesis includes
all smooth projective modifications, and their generic exceptional orders
recover point Lelong numbers by further point blowups. The source's
multiplier ideals use $e^{-t\varphi}$, as here.
Its Lelong numbers use growth against $\log\|z\|$, whereas our divisor
coefficient uses $\log|g|^2$; the source numbers are twice our numbers.
Equality of the two profiles is unchanged by this common factor.
The underlying local valuative theorem is \cite[Theorem A]{BFJ},
originally formulated for the stabilized ideals $\mathcal J_+$
with weight $e^{-2t\varphi}$. Strong openness
\cite[Theorem 1.1]{GZ} identifies $\mathcal J_+$ with $\mathcal J$;
rescaling $t$ then gives the convention used here.
For the reverse implication on an arbitrary smooth compact modification,
we record the usual local order estimate, avoiding an assumption that the
given modification is projective. Let $\varphi$ be a local potential, add
a smooth function to make it psh, and put
$b=\nu_E(\varphi\circ\pi)$ and
$e_k=\operatorname{ord}_E(\pi^{-1}\mathcal I(k\varphi)\cdot\mathcal O_Z)$.
Write $j_E\geq0$ for the generic order of the holomorphic Jacobian of
$\pi$. Then
\begin{equation}\label{eq:valuation-limit}
 b-\frac{j_E+1}{k}\leq\frac{e_k}{k}\leq b,
 \qquad\lim_{k\to\infty}\frac{e_k}{k}=b.
\end{equation}
Indeed, at a generic smooth point of $E=\{\zeta_1=0\}$, Siu subtraction
gives $\varphi\circ\pi\leq b\log|\zeta_1|^2+O(1)$.
Pulling back the integrability of any local germ $g\in\mathcal I(k\varphi)$
and retaining the Jacobian factor forces
$\operatorname{ord}_E(g\circ\pi)+j_E+1>kb$ by the one-variable radial
integral. This proves the lower estimate. For the upper estimate, the
local Bergman approximation \cite[Theorem 1.2]{DemB}, with its parameter
$m=k/2$, is
$\varphi_m=k^{-1}\log\sum_\ell|g_{m,\ell}|^2\geq\varphi-O(1/k)$.
All its holomorphic terms belong to $\mathcal I(k\varphi)$.
Local finite extraction of the Bergman ideal identifies its generic
order, giving
$e_k/k\leq\nu_E(\varphi_m\circ\pi)\leq b$.
Smooth additions affect neither ideal orders nor Lelong numbers.
Thus equality of all the original multiplier ideals implies equality of
the generic coefficients on every given modification. This completes
the equivalence; the local order estimate itself is established background.

Here is the envelope argument, with its known countable-divisor
domination step retained explicitly. The competitor $u$ shows that
$u\leq q\leq0$ and prevents collapse. The family is stable under
finite maxima, since maximum takes the minimum of the two generic
coefficients on every modification. Include $u$ in every finite
maximum of a Choquet sequence. The resulting increasing sequence
converges to $q$ almost everywhere. Its pullbacks are bounded below
by $u\circ\pi$ and above by zero; local $L^1$ convergence therefore
holds on each fixed modification. Upper semicontinuity of Lelong
numbers gives $\nu_E(q\circ\pi)\geq b(E)$, while $q\geq u$ gives
the reverse bound. Hence $q$ retains the exact ideal profile.

Take $C>A$ with $C\kappa\geq\theta$. The closed positive current
$T_q=C\kappa+\ddc q$ has coefficient $a_j$ along $H_j$.
Siu decomposition permits subtraction of every finite partial sum
of $S$; increasing these sums gives $T_q-S\geq0$. Consequently
\[
 (C-A)\kappa+\ddc(q-u)=T_q-S\geq0.
\]
The distributional difference has a genuine quasi-psh representative,
bounded above on compact $X$. Canonical uniqueness and almost-everywhere
order give $q\preceq u$ everywhere. Together with $q\geq u$ this
proves $q\sim u$. This mechanism is already known for countable
finite-mass divisorial potentials; it is not claimed as a new envelope
principle or a solution of the unrestricted zero-volume implications.

\section{Actual mass throughout the entire ordinary type}
Let $w\sim u$. It has the same base divisor coefficients, so
\[
 T=C\kappa+\ddc w\geq S,\qquad
 K=T-S=(C-A)\kappa+\ddc h\geq0,
 \qquad h=w-u.
\]
The canonical representative of the distribution $h$ is globally
bounded: its almost-everywhere bounds follow from ordinary equivalence,
and canonical uniqueness extends them. Thus $K$ has bounded local
potentials and its trace measure charges no pluripolar set. The
degree-one non-pluripolar current is precisely $\NP{T}=K$, since
$S$ is supported on the actual pole union and the remaining trace
measure charges none of it. In particular for every smooth
$\ddc$-closed $(1,1)$-form $\rho$,
\begin{equation}\label{eq:mixed}
 \int_X\rho\wedge\NP{T}=(C-A)\int_X\rho\wedge\kappa.
\end{equation}
This is distributional Stokes with the bounded function $h$.

To compute the closed top mass, first use $u$ itself. Its finite
partial sums $u_N=\sum_{j\leq N}a_j\ell_j$, up to harmless constants,
have closed actual mass $(C-\sum_{j\leq N}a_j)^2$: away from the
finite divisor union their closed current is the smooth residual.
Since $u_0\leq u_N$, mass monotonicity in the fixed closed K\"ahler
background $C\kappa$ \cite[Theorem 1.2]{WN} gives
\[
 \int_X\NP{(C\kappa+\ddc u)^2}
 \leq\left(C-\sum_{j\leq N}a_j\right)^2.
\]
Conversely, positivity and multilinearity of closed non-pluripolar
products give a lower bound $(C-A)^2$ from the summand
$(C-A)\kappa$ of $C\kappa+\ddc u=S+(C-A)\kappa$
\cite[Propositions 1.4(c), 1.6]{BEGZ}. Letting $N\to\infty$ gives
equality. Since $w\sim u$, applying the same monotonicity theorem
in both directions gives
\begin{equation}\label{eq:closed}
 \int_X\NP{T^2}=(C-A)^2.
\end{equation}
The compact manifold and every current in this argument are K\"ahler
and closed respectively. No nonclosed mass-monotonicity theorem is used.

Put $L=C\kappa-\theta$, so $\ddc L=0$. At bounded truncation levels
the smooth-background expansion, restricted to the strict superlevel,
gives the actual Borel measure identity
\begin{equation}\label{eq:correction}
 \operatorname{MA}^{\mathrm{ext}}_\theta(w)
 =\NP{T^2}-2L\wedge\NP{T}+L^2.
\end{equation}
Indeed use the same truncations $\max(w,-j)$ in both backgrounds,
expand the bounded products, and restrict to $\{w>-j\}$.
The closed products are finite, their degree-one limit is $K$, and
the smooth correction charges no pole set. The limit therefore gives
\eqref{eq:correction} with finite signed correction measures.
This is the bounded local correction mechanism of
\cite[Section 2.1]{BGL}; a general unbounded theorem requiring BMP
is not a substitute for this explicit passage.

Integrating \eqref{eq:correction}, using \eqref{eq:mixed} with
$\rho=L$ and \eqref{eq:closed}, gives
\[
 \int_X\operatorname{MA}^{\mathrm{ext}}_\theta(w)
 =(C-A)^2-2(C-A)\int_XL\wedge\kappa+\int_XL^2
 =\int_X(\theta-A\kappa)^2.
\]
This proves \eqref{eq:mass} for every representative, rather than
identifying one envelope's raw mass with an extremal volume.

\section{A concrete nonclosed and nonanalytic example}
In a coordinate ball with coordinates $z_1,z_2$, choose a smooth
real cutoff $\chi$ equal to one near the origin. Set
$\alpha=\chi|z_2|^2dz_1$, extended by zero, and
\[
 \tau=\partial\overline\alpha+\bar\partial\alpha.
\]
This is smooth and real, with $\ddc\tau=0$. Where $\chi=1$,
\[
 d\tau=dz_2\wedge d\bar z_2\wedge(dz_1-d\bar z_1)\ne0.
\]
For small $\epsilon>0$ put
$\theta=\kappa+\epsilon\tau\geq\kappa/2$.
Take $A=1/4$, $a_j=A2^{-j}$, $t_j=2^{-j}$ and
\[
 H_j=\{Z_0+t_jZ_1+t_j^2Z_2=0\}.
\]
Any three defining coefficient vectors are independent by the
Vandermonde determinant. The divisors are therefore outside a
single pencil, and all have positive prescribed coefficients.
An analytic ideal type on a compact base has only finitely many
prime divisorial components, so this ideal type is nonanalytic.
Here $R\geq\kappa/4$, and the two singularity volumes are at least
$1/16$. The background remains genuinely nonclosed and pluriclosed.

For any smooth gauge $\theta'=\theta+\ddc f$, pair a potential with
$w'=w-f$. The actual current is unchanged, ideals are unchanged and
ordinary comparisons are preserved. This bijection of entire ordinary
types preserves both singularity volumes. It does not identify a
normalized zero-obstacle envelope with $q-f$ without projection.
No unrestricted modification law, zero-volume converse or general
nonclosed ideal classification is asserted.

\begin{thebibliography}{9}
\bibitem{Dem} J.-P. Demailly, \emph{Complex Analytic and Differential
Geometry}, June 2012 manuscript: canonical representatives,
Poincar\'e--Lelong formula and Siu decomposition.
\bibitem{DX} T. Darvas and M. Xia, \emph{The closures of test
configurations and algebraic singularity types}, Adv. Math. 397 (2022),
108198, Theorem 2.13 and Proposition 2.30.
\bibitem{BFJ} S. Boucksom, C. Favre and M. Jonsson,
\emph{Valuations and plurisubharmonic singularities},
arXiv:math/0702487v2, Theorem A.
\bibitem{GZ} Q. Guan and X. Zhou,
\emph{A proof of Demailly's strong openness conjecture},
Ann. Math. 182 (2015), 605--616, Theorem 1.1.
\bibitem{DemB} J.-P. Demailly, \emph{On the cohomology of pseudoeffective
line bundles}, arXiv:1401.5432v2, Theorem 1.2 and Lemma 1.10.
\bibitem{BT} E. Bedford and B. A. Taylor, \emph{A new capacity for
plurisubharmonic functions}, Acta Math. 149 (1982), 1--40.
\bibitem{Gau} P. Gauduchon, \emph{Les m\'etriques standard d'une surface
complexe compacte \`a premier nombre de Betti pair}, Ast\'erisque 126
(1985), 129--135, Theorem 2.
\bibitem{BEGZ} S. Boucksom, P. Eyssidieux, V. Guedj and A. Zeriahi,
\emph{Monge--Amp\`ere equations in big cohomology classes},
Acta Math. 205 (2010), 199--262.
\bibitem{WN} D. Witt Nystr\"om, \emph{Monotonicity of non-pluripolar
Monge--Amp\`ere masses}, arXiv:1703.01950v1.
\bibitem{BGL} S. Boucksom, V. Guedj and C. H. Lu,
\emph{Volumes of Bott--Chern classes}, arXiv:2406.01090, Section 2.1.
\end{thebibliography}
\end{document}
